\maketitle
\makesignature

\ifproject
\begin{abstractTH}
โครงงานนี้ประเมินแนวทางการสร้าง Investigation Context จาก Security Logs ในระบบ SIEM
ก่อนส่งให้โมเดลภาษาขนาดใหญ่ (LLM) วิเคราะห์ เพื่อแก้ปัญหาข้อมูลจำนวนมาก, Noise
และต้นทุน Token ที่สูงจากการใช้ Raw Logs โดยตรง
โครงงานพัฒนา Context Construction Pipeline ที่คัดกรองและจัดรูปแบบเฉพาะข้อมูลที่เกี่ยวข้องกับเหตุการณ์
ร่วมกับการอ้างอิงองค์ความรู้จาก MITRE ATT\&CK เพื่อให้ LLM ขนาดเล็กที่ใช้ทำงานภายในเครื่อง
สามารถระบุเทคนิคการโจมตีได้อย่างแม่นยำ
ผลลัพธ์จะถูกนำเสนอเป็นรายงานการสืบสวน (Investigation Report) ผ่าน Web Dashboard
เพื่อช่วยให้ SOC Analyst ลดเวลาในการสืบสวนเหตุการณ์ และลดต้นทุนในการใช้ LLM
โดยยังคงหลักฐานที่จำเป็นต่อการวิเคราะห์ไว้ครบถ้วน
\end{abstractTH}

\begin{abstract}
This project evaluates approaches for constructing an investigation context from
security logs collected by a Security Information and Event Management (SIEM) system
before they are analyzed by a large language model (LLM).
Feeding raw logs directly to an LLM is problematic because of the large data volume,
the amount of irrelevant information (noise), and the resulting high token cost.
We develop a Context Construction Pipeline that normalizes the logs, selects only the
events related to the alerted incident, and formats them in a compact representation,
combined with knowledge retrieved from the MITRE ATT\&CK framework, so that a small
LLM running locally can accurately identify the attack technique.
The analysis is presented as an investigation report on a web dashboard, helping SOC
analysts reduce investigation time and LLM usage cost while preserving the evidence
required for the analysis.
\end{abstract}

\iffalse
\begin{dedication}
Dedication page is optional.
\end{dedication}
\fi % \iffalse

% \begin{acknowledgments}
% โครงงานนี้สำเร็จลุล่วงได้ด้วยความกรุณาจาก ผศ.ดร.ยุทธพงษ์ สมจิต อาจารย์ที่ปรึกษาโครงงาน
% ที่ได้สละเวลาให้คำปรึกษา ข้อเสนอแนะ และแนวทางในการดำเนินงานตลอดระยะเวลาการทำโครงงาน
% คณะผู้จัดทำขอกราบขอบพระคุณเป็นอย่างสูงมา ณ ที่นี้

% ขอขอบพระคุณคณะกรรมการสอบโครงงาน ได้แก่ รศ.ดร.ตรัสพงศ์ ไทยอุปถัมภ์
% ผศ.ดร.ศศิน จันทร์พวงทอง และ ผศ.ดร.กานต์ ปทานุคม
% ที่ได้ให้คำถามและข้อเสนอแนะอันเป็นประโยชน์ระหว่างการนำเสนอ
% ซึ่งช่วยให้คณะผู้จัดทำปรับปรุงการนิยามปัญหา ขอบเขต และวิธีการทดลองให้ชัดเจนยิ่งขึ้น

% ขอขอบคุณภาควิชาวิศวกรรมคอมพิวเตอร์ คณะวิศวกรรมศาสตร์ มหาวิทยาลัยเชียงใหม่
% ที่ให้ความรู้และสนับสนุนการทำโครงงาน
% รวมถึงผู้พัฒนาชุดข้อมูลสาธารณะ Splunk Attack Data และ OTRF Security-Datasets
% ที่ทำให้การทดลองในโครงงานนี้เป็นไปได้

% \acksign{2026}{10}{22}
% \end{acknowledgments}%
\fi % \ifproject

\contentspage

\ifproject
\figurelistpage

\tablelistpage
\fi % \ifproject
